\documentclass{article}

% ====== Paquetes básicos ======
\usepackage{tikz}               % Diagramas
\usetikzlibrary{babel,decorations.pathmorphing}
\usepackage[utf8]{inputenc}    % Codificación UTF-8
\usepackage[T1]{fontenc}       % Acentos correctos
\usepackage[spanish]{babel}    % Traducción al español
\usepackage{csquotes}          % Recomendado para biblatex con babel
\usepackage{hyperref}          % Hipervínculos
\usepackage{amsmath,amssymb}   % Matemáticas
\usepackage{graphicx}          % Imágenes
\usepackage{caption}           % Mejor control de captions
\usepackage{geometry}          % Márgenes
\usepackage{fancyhdr}          % Encabezados y pies
\usepackage{parskip}           % Mejor separación entre párrafos

% ====== Referencias ======
\usepackage[
    backend=biber,
    style=numeric-comp,     % cambie a APA de ser necesario
    sorting=none
]{biblatex} 
\bibliography{references}

% agrege `\printbibliography{}` para las referencias
% cite con `\cite{}`
% agregue las referencias a un `references.bib`
% citar todo con
% \nocite{*}
% \printbibliography{}

% ====== Configuración ======
\geometry{margin=2.5cm}
\pagestyle{fancy}
\fancyhf{}
\rhead{\thepage}
\setlength{\headheight}{15pt}
\lhead{\leftmark}

% ====== Comando personalizado para imágenes ======
% Uso: \imagen{ruta}{Texto del pie de figura}
\newcommand{\imagen}[2]{
    \begin{figure}[h!]
        \centering
        \includegraphics[width=200px]{#1}
        \caption{#2}
    \end{figure}
}

% ====== Datos del artículo ======
\title{Portada}
\author{Lester Iván Hernández Hernandez}


\usepackage[spanish]{babel}
\usepackage{datetime}
\newdateformat{mydateformat}{\THEDAY/\THEMONTH/\THEYEAR}

\newcommand{\titlename}{Memoria académica del Congreso de Ciencias}
\newcommand{\signature}{Procesos Estocásticos (MM-435)}
\newcommand{\seccion}{1200}
\newcommand{\docente}{Ph.D. Marcos Orellana}
%Individual
\newcommand{\alumno}{Lester Iván Hernández Hernandez}
\newcommand{\cuenta}{20242000637}

%Grupo
%\newcommand{\grupo}{2}
%\newcommand{\integrantes}{}

\newcommand{\portada}{
    \begin{titlepage}
        \centering
        \vspace*{-1cm}
        \includegraphics[width=4.5cm]{Logo-unah.jpg}\\
        \vspace*{0.5cm}
        {\large\bfseries UNIVERSIDAD NACIONAL AUTÓNOMA DE HONDURAS EN EL VALLE DE SULA
        \\UNAH-VS\par}
        \vspace{2cm}
        {\large \titlename\par}
        \vspace{1cm}
        {\large Asignatura: \signature\par}
        {\large Sección: \seccion\par}
        {\large Docente: \docente\par}
        \vspace{1cm}
        %Individual
        {\large Alumno: \alumno\par}
        {\large N. de cuenta: \cuenta\par}
        %Grupo
        %{\large Grupo: \grupo\par}
        %{\large Integrantes: \integrantes\par}
        \vfill
        {\large \today\par}
        \vspace*{1in}
    \end{titlepage}
}


\begin{document}
    \portada{}
    \newpage

    \section{Introducción}
    \textit{Presente brevemente el Congreso de Ciencias, indique el propósito de esta memoria y explique el criterio utilizado para seleccionar las conferencias o cursos incluidos.}

    \section{Modelos Matemáticos Discretos}
    \textbf{Expositor:} Miguel Piñar

    \subsection{Descripción general}
    \textit{Describa el contexto de la conferencia, el tema abordado, su importancia y los principales aspectos desarrollados.}

    \subsection{Objetivo principal}
    \textit{Explique el propósito central de la conferencia.}

    \subsection{Principales hallazgos o resultados}
    \textit{Describa los resultados, conclusiones, aportaciones o conocimientos más relevantes presentados.}

    \subsection{Reflexión personal}
    \textit{Explique qué conocimientos adquirió y cuál es su importancia o aplicación en su formación académica o profesional.}

    \section{Las matemáticas tras el éxito de Google}
    \textit{Título alternativo: Cinco cosas que Google no quiere que sepas}

    \textbf{Expositor:} Miguel Piñar

    \subsection{Descripción general}
    \textit{Describa el contexto de la conferencia, el tema abordado, su importancia y los principales aspectos desarrollados.}

    \subsection{Objetivo principal}
    \textit{Explique el propósito central de la conferencia.}

    \subsection{Principales hallazgos o resultados}
    \textit{Describa los resultados, conclusiones, aportaciones o conocimientos más relevantes presentados.}

    \subsection{Reflexión personal}
    \textit{Explique qué conocimientos adquirió y cuál es su importancia o aplicación en su formación académica o profesional.}

    \section{Impacto del taller ``Prepárate para la PHUMA'' en el aprendizaje matemático}
    \textbf{Título completo:} Impacto del taller ``Prepárate para la PHUMA'' en el aprendizaje matemático de los estudiantes del Instituto Perla del Ulúa: un análisis mediante indicadores estadísticos

    \textbf{Expositores:} Carmen Flores, Omar Pineda y Lino Ávila

    \subsection{Descripción general}
    \textit{Describa el contexto de la investigación, el taller desarrollado, la población estudiada y los principales aspectos metodológicos presentados.}

    \subsection{Objetivo principal}
    \textit{Explique el propósito central de la investigación o presentación.}

    \subsection{Principales hallazgos o resultados}
    \textit{Describa los indicadores estadísticos, resultados, conclusiones y aportaciones más relevantes.}

    \subsection{Reflexión personal}
    \textit{Explique qué conocimientos adquirió y cuál es la importancia o aplicación de lo aprendido en su formación académica o profesional.}

    \section{Modelado teórico y simulación estadística de variables biológicas y de salud}
    \textbf{Título completo:} Modelado teórico y simulación estadística de variables biológicas y de salud mediante distribuciones de probabilidad asimétricas inspiradas en la Física

    \textbf{Expositor:} Marcos Orellana

    \subsection{Descripción general}
    \textit{Describa el problema abordado, los modelos teóricos, las distribuciones de probabilidad y las simulaciones presentadas.}

    \subsection{Objetivo principal}
    \textit{Explique el propósito central de la conferencia.}

    \subsection{Principales hallazgos o resultados}
    \textit{Describa los resultados, conclusiones, aportaciones o conocimientos más relevantes presentados.}

    \subsection{Reflexión personal}
    \textit{Explique qué conocimientos adquirió y cuál es su importancia o aplicación en su formación académica o profesional.}

    \section{Gemas de la Física teórica}
    \textit{Título alternativo: El código secreto del universo: la raíz cuadrada de la gravedad, de las ecuaciones de Maxwell a la doble copia}

    \textbf{Expositor:} Bryan Larios

    \subsection{Descripción general}
    \textit{Describa el contexto de la conferencia, los conceptos de Física teórica abordados y los principales aspectos desarrollados.}

    \subsection{Objetivo principal}
    \textit{Explique el propósito central de la conferencia.}

    \subsection{Principales hallazgos o resultados}
    \textit{Describa los resultados, conclusiones, aportaciones o conocimientos más relevantes presentados.}

    \subsection{Reflexión personal}
    \textit{Explique qué conocimientos adquirió y cuál es su importancia o aplicación en su formación académica o profesional.}

    \section{Conclusiones generales}
    \textit{Integre las principales conclusiones de la memoria y explique cómo las conferencias seleccionadas contribuyeron a su formación académica o profesional.}


\end{document} 
